% p.61 — الباب السابع والعشرون في المذكر والمؤنث والمفرد والجمع / الباب الثامن والعشرون في كل الافعال المستعملة
\mspage{61}
% Layout (as in the lithograph): the page is split vertically into two boxed tables printed
% side by side in one `tables` row, each carrying its chapter heading as a \multicolumn row
% at its top (there is no page-wide heading).
% The RIGHT (wider) box: heading of bab 27 and a four-column table (مذكر | مؤنث | مفرد | جمع).
% The two right-hand columns (مذكر/مؤنث) are subdivided by two internal sub-header rows
% (مفرد | تثنية, then مفرد | جمع), while the two left-hand columns (مفرد | جمع) run as one
% continuous list of 20 rows. The sub-header rows are reproduced as \thead cells in the rows
% where they fall on the page.
% The LEFT (narrower) box: heading of bab 28 (two printed lines, two \multicolumn rows) and a
% two-column list of the verb patterns (رقم + وزن), with group headings inside the columns
% (narrow p-columns, so the long group headings wrap as they do on the page).
\begin{tables}
\begin{sarf}{|c|c|c|c|}
\multicolumn{4}{|c|}{\thead{الباب السابع والعشرون في المذكر والمؤنث والمفرد والجمع}}\\\hline
\thead{مذكر} & \thead{مؤنث} & \thead{مفرد} & \thead{جمع}\\\hline
ضَارِبٌ & ضَارِبَةٌ & غُلْبَةٌ & غُلْبٌ\\
عَطْشَانُ & عَطْشٰى & صُوْرَةٌ & صُوَرٌ\\
عَطْشَانٌ & عَطْشَانَةٌ & حِلْيَةٌ & حِلًى\\
أَبْيَضُ & بَيْضَاءُ & بُلْبُلٌ & بَلَابِلُ\\
أَكْرَمُ & كُرْمٰى & أُسْلُوْبٌ & أَسَالِيْبُ\\
طَالِقٌ & طَالِقٌ & قِرْطَاسٌ & قَرَاطِيْسُ\\
عَلَّامَةٌ & عَلَّامَةٌ & مُسْلِمٌ & مُسْلِمُوْنَ\\\cline{1-2}
\thead{مفرد} & \thead{تثنية} & مُسْلِمَةٌ & مُسْلِمَاتٌ\\\cline{1-2}
رَجُلٌ & رَجُلَانِ & رَجُلٌ & رِجَالٌ\\
ذُوْ مَالٍ & ذَوَا مَالٍ & رَسُوْلٌ & رُسُلٌ\\
ذِكْرٰى & ذِكْرَيَانِ & رَاكِبٌ & رَكْبٌ\\
عَصَا & عَصَوَانِ & عَاذِلٌ & عَوَاذِلُ\\
قَهْقَرٰى & قَهْقَرَوَانِ & غَنِيٌّ & أَغْنِيَاءُ\\
أَبٌ & أَبَوَانِ & كُرْمٰى & كُرْمَيَاتٌ\\
رَامٍ & رَامِيَانِ & أَعْظَمُ & أَعْظَمُوْنَ\\\cline{1-2}
\thead{مفرد} & \thead{جمع} & أَبْيَضُ & بِيْضٌ\\\cline{1-2}
طِفْلٌ & أَطْفَالٌ & دَاعٍ & دُعَاةٌ\\
نَهْرٌ & أَنْهُرٌ & سَاجِدٌ & سُجَّدًا\\
\com{رَغِيْقٌ}{كذا في الأصل، ولعل الصواب: رَغِيْفٌ} & أَرْغِفَةٌ & مَرِيْضٌ & مَرْضٰى\\
فَتًى & فِتْيَةٌ & سَكْرَانُ & سُكَارٰى\\
\end{sarf}
% TODO: the chapter number printed after «والعشرون» in the heading of bab 28 reads like ۲۱;
% expected ۲۸.  Transcribed as \unclear{21}.
% Left box: heading of bab 28, then two columns, read right column first (top to bottom), then
% left column. Group headings (ثلاثي مجرد etc.) sit inside the columns as cells.
\begin{sarf}{|>{\centering\arraybackslash}p{22mm}|>{\centering\arraybackslash}p{22mm}|}
\multicolumn{2}{|c|}{\thead{الباب الثامن والعشرون \unclear{21}}}\\
\multicolumn{2}{|c|}{\thead{في كل الافعال المستعملة}}\\\hline
\thead{ثلاثي مجرد} & (15) اِسْتَفْعَلَ\\
(1) فَتَحَ & (16) اِفَّعَّلَ\\
(2) ضَرَبَ & (17) اِفَّاعَلَ\\
(3) نَصَرَ & (18) اِفْعَالَّ\\
(4) سَمِعَ & (19) اِفْعَوَّلَ\\
(5) حَسِبَ & (20) اِفْعَوْعَلَ\\
(6) كَرُمَ & \thead{رباعي مجرد}\\
\thead{ثلاثي مزيد بحرف} & (21) فَعْلَلَ\\
(7) أَفْعَلَ & \thead{رباعي مزيد بحرف}\\
(8) فَعَّلَ & (22) تَفَعْلَلَ\\
(9) فَاعَلَ & \thead{رباعي مزيد بحرفين}\\
\thead{ثلاثي مزيد بحرفين} & (23) اِفْعَنْلَلَ\\
(10) تَفَعَّلَ & (24) اِفْعَلَلَّ\\
(11) تَفَاعَلَ & \thead{ملحق بالرباعي المجرد دَحْرَجَ}\\
(12) اِفْتَعَلَ & (1) جَلْبَبَ\\
(13) اِنْفَعَلَ & (2) قَلْنَسَ\\
(14) اِفْعَلَّ & (3) جَوْرَبَ\\
\thead{ثلاثي مزيد بثلاثة احرف} & (4) سَرْوَلَ\\
\end{sarf}
\end{tables}
